\documentclass[12pt,a4paper]{article}

% ============================================================
% PACKAGES
% ============================================================
\usepackage[utf8]{inputenc}
\usepackage[T1]{fontenc}
\usepackage{lmodern}
\usepackage[french]{babel}
\usepackage{geometry}
\usepackage{graphicx}
\usepackage[table]{xcolor}
\usepackage{listings}
\usepackage{booktabs}
\usepackage{array}
\usepackage{fancyhdr}
\usepackage{titlesec}
\usepackage{enumitem}
\usepackage{mdframed}
\usepackage{float}
\usepackage{microtype}
\usepackage{setspace}
\usepackage{amssymb}
\usepackage{hyperref}
\setlength{\emergencystretch}{2em}

% ============================================================
% MISE EN PAGE
% ============================================================
\geometry{
  headheight=25pt,
  top=2.6cm,
  bottom=2.8cm,
  left=2.8cm,
  right=2.8cm
}
\setstretch{1.12}

% ============================================================
% COULEURS
% ============================================================
\definecolor{primary}{RGB}{28, 90, 168}
\definecolor{primarylight}{RGB}{88, 145, 220}
\definecolor{secondary}{RGB}{232, 240, 252}
\definecolor{codebg}{RGB}{248, 250, 255}
\definecolor{codeframe}{RGB}{170, 195, 230}
\definecolor{success}{RGB}{34, 130, 90}
\definecolor{warning}{RGB}{200, 120, 40}
\definecolor{danger}{RGB}{185, 40, 40}
\definecolor{darkgray}{RGB}{45, 55, 72}
\definecolor{lightgray}{RGB}{95, 110, 140}
\definecolor{rowgray}{RGB}{245, 247, 250}

% ============================================================
% STYLE DU CODE
% ============================================================
\lstset{
  backgroundcolor=\color{codebg},
  frame=single,
  rulecolor=\color{codeframe},
  basicstyle=\ttfamily\footnotesize,
  keywordstyle=\color{primary}\bfseries,
  commentstyle=\color{lightgray}\itshape,
  stringstyle=\color{success},
  breaklines=true,
  showstringspaces=false,
  numbers=left,
  numberstyle=\tiny\sffamily\textcolor{lightgray},
  numbersep=6pt,
  tabsize=2,
  captionpos=b,
  aboveskip=1em,
  belowskip=1em,
  xleftmargin=0pt,
  framexleftmargin=0pt
}

% ============================================================
% EN-TETE ET PIED DE PAGE
% ============================================================
\pagestyle{fancy}
\fancyhf{}
\fancyhead[L]{\small\sffamily\textcolor{primary!85!black}{Projet M1 GIL --- Gestion de flotte}}
\fancyhead[R]{\small\sffamily\textcolor{primary!85!black}{Univ. Rouen Normandie \textperiodcentered{} 2025--2026}}
\fancyfoot[C]{\small\sffamily\textcolor{primary!70!black}{\thepage}}
\renewcommand{\headrulewidth}{0.4pt}
\renewcommand{\headrule}{\hbox to\headwidth{\color{primary!35}\leaders\hrule height \headrulewidth\hfill}}
\renewcommand{\footrulewidth}{0pt}

% ============================================================
% STYLE DES TITRES
% ============================================================
\titleformat{\section}
  {\Large\sffamily\bfseries\color{primary}}
  {\thesection.}{0.55em}{}

\titleformat{\subsection}
  {\large\sffamily\bfseries\color{primary!85!black}}
  {\thesubsection.}{0.45em}{}

\titleformat{\subsubsection}
  {\normalsize\sffamily\bfseries\color{primary!75!black}}
  {\thesubsubsection.}{0.4em}{}

\titlespacing*{\section}{0pt}{1.6ex plus .2ex}{1ex plus .1ex}
\titlespacing*{\subsection}{0pt}{1.3ex plus .15ex}{0.7ex plus .1ex}

% ============================================================
% BOITE INFO
% ============================================================
\newmdenv[
  backgroundcolor=secondary,
  linecolor=primary,
  linewidth=1.1pt,
  leftline=true,
  rightline=false,
  topline=false,
  bottomline=false,
  innerleftmargin=11pt,
  innerrightmargin=11pt,
  innertopmargin=9pt,
  innerbottommargin=9pt,
  skipabove=0.6em,
  skipbelow=0.6em
]{infobox}

\hypersetup{
  colorlinks=true,
  linkcolor=primary!90!black,
  urlcolor=primarylight,
  pdfborderstyle={/S/U/W 0}
}

% ============================================================
% DOCUMENT
% ============================================================
\begin{document}

% ------------------------------------------------------------
% PAGE DE TITRE
% ------------------------------------------------------------
\begin{titlepage}
  \centering
  \vspace*{1cm}

  {\Large\textbf{Universit\'{e} de Rouen Normandie}\\[0.3cm]
  \large Master 1 -- G\'{e}nie Informatique et Logiciel}

  \vspace{1.5cm}
  {\color{primary!55}\rule{\linewidth}{0.5mm}}\\[0.45cm]

  {\Huge\sffamily\bfseries\color{primary} Projet Gestion de Flotte\\[0.35cm]}
  {\LARGE\sffamily\bfseries\color{primary!80!black} Rapport de Semaine 7}\\[0.35cm]
  {\large\sffamily\color{primary!65!black}
    Observabilit\'{e} compl\`{e}te \textperiodcentered{}
    Alertes temps r\'{e}el \textperiodcentered{}
    Tests de charge K6 \textperiodcentered{}
    D\'{e}ploiement Helm prod-ready}

  {\color{primary!55}\rule{\linewidth}{0.5mm}}\\[1.5cm]

  \begin{minipage}{0.45\textwidth}
    \begin{flushleft}
      \textbf{\'{E}tudiants :}\\
      Ahc\`{e}ne Amouchas \\
      Florian Pepin \\[0.5cm]
      \textbf{Encadrants :}\\
      Lydia \& Luc
    \end{flushleft}
  \end{minipage}
  \hfill
  \begin{minipage}{0.45\textwidth}
    \begin{flushright}
      \textbf{Ann\'{e}e universitaire :}\\
      2025 -- 2026 \\[0.5cm]
      \textbf{Semaine :}\\
      7 / 7
    \end{flushright}
  \end{minipage}

  \vfill
  {\large 23 avril 2026}
\end{titlepage}

\tableofcontents
\newpage

% ============================================================
\section{Observabilit\'{e} compl\`{e}te}
% ============================================================

\subsection{Architecture de la stack}

L'observabilit\'{e} de la plateforme repose sur le pattern \textbf{LGTMo}
(Loki, Grafana, Tempo/Jaeger, M\'{e}triques via Prometheus, OpenTelemetry),
int\'{e}gralement pilot\'{e} par le chart Helm \texttt{fleet-observability}.
L'\textbf{OpenTelemetry Collector} constitue le pivot central : il re\c{c}oit
traces, m\'{e}triques et logs de tous les microservices via OTLP et les route
vers les backends sp\'{e}cialis\'{e}s.

\begin{lstlisting}[language=bash, caption=Pipeline de l'OTel Collector
--- \texttt{otel-collector-config.yaml}]
receivers:
  otlp:
    protocols:
      grpc:  { endpoint: "0.0.0.0:4317" }
      http:  { endpoint: "0.0.0.0:4318" }

processors:
  memory_limiter: { limit_percentage: 70, spike_limit_percentage: 30 }
  batch: {}

exporters:
  otlp/jaeger: { endpoint: "jaeger:4317", tls: { insecure: true } }
  loki:        { endpoint: "http://loki:3100/loki/api/v1/push" }
  prometheus:  { endpoint: "0.0.0.0:8889" }

service:
  pipelines:
    traces:  { receivers: [otlp], processors: [batch], exporters: [otlp/jaeger] }
    metrics: { receivers: [otlp], processors: [batch], exporters: [prometheus] }
    logs:    { receivers: [otlp], processors: [batch], exporters: [loki] }
\end{lstlisting}

\subsection{Jaeger --- Distributed Tracing}

Jaeger (\texttt{jaegertracing/all-in-one:1.55}) est d\'{e}ploy\'{e} en mode
\textit{all-in-one} pour la d\'{e}monstration. Il re\c{c}oit les traces via
OTLP gRPC (port 4317) depuis l'OTel Collector et expose son interface
graphique sous le chemin \texttt{/jaeger} de l'ingress \texttt{flotte.local},
ce qui permet une navigation unifi\'{e}e depuis le navigateur.

\begin{infobox}
\textbf{Correction technique appliqu\'{e}e :} La datasource Jaeger dans Grafana
doit pointer vers \texttt{jaeger:16686\textbf{/jaeger}} (avec le base path) et
non \texttt{jaeger:16686}. Sans ce correctif, l'exploration de traces
depuis les panels Grafana retourne une erreur 404.
\end{infobox}

\subsection{Prometheus --- M\'{e}triques}

Prometheus (\texttt{prom/prometheus:v2.51.0}) scrape les m\'{e}triques agr\'{e}g\'{e}es
par l'OTel Collector sur le port 8889 avec un intervalle de \textbf{5 secondes},
adapt\'{e} au suivi en quasi-temps r\'{e}el des \'{e}v\'{e}nements de la flotte. Les
microservices ne s'exposent pas directement \`{a} Prometheus~: ils poussent
leurs m\'{e}triques en OTLP, ce qui centralise la configuration de collecte.

\begin{lstlisting}[language=bash, caption=Configuration du scraping Prometheus]
global:
  scrape_interval: 5s

scrape_configs:
  - job_name: 'otel-collector'
    static_configs:
      - targets: ['otel-collector:8889']
\end{lstlisting}

\subsection{Loki + Promtail --- Centralisation des logs}

Loki (\texttt{grafana/loki:2.9.4}) centralise les logs applicatifs. Promtail
est d\'{e}ploy\'{e} en \textbf{DaemonSet} pour collecter les logs au niveau de chaque
n\oe{}ud Kubernetes. Cette approche compl\`{e}te l'OTel Collector pour les
composants qui n'\'{e}mettent pas nativement en OTLP (frontend React,
\texttt{location-service} Node.js).

\begin{infobox}
\textbf{Probl\`{e}me rencontr\'{e} sous Minikube :} Promtail ne trouvait pas les
fichiers de log \`{a} cause de symlinks Docker absents. Solution appliqu\'{e}e~:
configuration \texttt{static\_configs} pointant directement vers
\texttt{/var/log/pods/**/*.log} au lieu des symlinks \texttt{/var/log/containers}.
\end{infobox}

\subsection{Alertes m\'{e}tier}

Les alertes m\'{e}tier sont produites par le \texttt{events-service} \`{a} partir
de la consommation des topics Kafka. La cha\^{i}ne compl\`{e}te est la suivante~:

\begin{enumerate}[leftmargin=*]
  \item Un \'{e}v\'{e}nement arrive sur un topic Kafka m\'{e}tier
        (\texttt{location.events}, \texttt{alertes.events} ou
        \texttt{conducteurs.events}).
  \item Le \texttt{events-service} consomme l'\'{e}v\'{e}nement, cr\'{e}e une alerte
        en base et publie \texttt{ALERT\_CREATED} sur
        \texttt{flotte.alerts.created}.
  \item La Gateway GraphQL propage l'alerte via la subscription WebSocket
        \texttt{ALERT\_CREATED\_SUB} vers tous les clients connect\'{e}s.
  \item Le frontend met \`{a} jour la vue \texttt{/alerts} sans rechargement.
\end{enumerate}

\begin{table}[H]
\centering
\small
\renewcommand{\arraystretch}{1.4}
\begin{tabular}{
  >{\ttfamily\raggedright\arraybackslash}p{4.8cm}
  >{\ttfamily\raggedright\arraybackslash}p{5cm}
  >{\centering\arraybackslash}p{3cm}
}
\toprule
\rowcolor{secondary}
\normalfont\textbf{Type d'alerte} & \normalfont\textbf{Topic source} & \normalfont\textbf{S\'{e}v\'{e}rit\'{e}} \\
\midrule
SPEED\_EXCEEDED        & flotte.location.events    & WARNING / HIGH \\
\rowcolor{rowgray}
GEOFENCING\_BREACH     & flotte.location.events    & HIGH \\
VEHICLE\_IMMOBILIZED   & flotte.location.events    & WARNING \\
\rowcolor{rowgray}
MAINTENANCE\_OVERDUE   & flotte.alertes.events     & CRITICAL / HIGH \\
LICENSE\_EXPIRING      & flotte.conducteurs.events & WARNING \\
\rowcolor{rowgray}
LICENSE\_EXPIRED       & flotte.conducteurs.events & CRITICAL \\
\bottomrule
\end{tabular}
\caption{Types d'alertes m\'{e}tier et topics Kafka associ\'{e}s}
\end{table}

% ============================================================
\section{Dashboards Grafana personnalis\'{e}s}
% ============================================================

\subsection{Configuration des datasources}

Grafana 10.4.1 est configur\'{e} avec trois datasources d\'{e}clar\'{e}es en tant que
\textit{provisioned resources} (fichiers YAML mont\'{e}s dans le pod), ce qui
\'{e}vite toute configuration manuelle apr\`{e}s d\'{e}ploiement.

\begin{table}[H]
\centering
\small
\renewcommand{\arraystretch}{1.4}
\begin{tabular}{
  >{\raggedright\arraybackslash}p{3.2cm}
  >{\ttfamily\raggedright\arraybackslash}p{6cm}
  >{\raggedright\arraybackslash}p{4.5cm}
}
\toprule
\rowcolor{secondary}
\normalfont\textbf{Datasource} & \normalfont\textbf{URL interne} & \normalfont\textbf{Usage} \\
\midrule
Prometheus (d\'{e}faut) & http://prometheus:9090    & M\'{e}triques techniques et m\'{e}tier \\
\rowcolor{rowgray}
Loki                & http://loki:3100           & Logs applicatifs \\
Jaeger              & http://jaeger:16686/jaeger & Traces distribu\'{e}es \\
\bottomrule
\end{tabular}
\caption{Datasources Grafana provisionn\'{e}es automatiquement}
\end{table}

\subsection{Dashboards m\'{e}tier (business metrics)}

Les dashboards couvrent quatre domaines fonctionnels~:

\begin{itemize}[leftmargin=*]
  \item \textbf{Vue flotte} : nombre de v\'{e}hicules disponibles, en maintenance,
        en tourn\'{e}e, hors service. \'{E}volution temporelle du statut du parc.
  \item \textbf{Vue conducteurs} : conducteurs actifs, en tourn\'{e}e, indisponibles.
        Alertes permis expirant sous 30 jours.
  \item \textbf{Vue alertes} : r\'{e}partition des alertes actives par s\'{e}v\'{e}rit\'{e}
        (CRITICAL, HIGH, WARNING, INFO). Taux d'acquittement.
  \item \textbf{Vue performance API} : latences p50/p95/p99 des endpoints
        GraphQL, REST et gRPC. Taux d'erreur par service.
\end{itemize}

Grafana est accessible sous \texttt{http://flotte.local/grafana} (chemin
relatif configur\'{e} via \texttt{GF\_SERVER\_ROOT\_URL} et
\texttt{GF\_SERVER\_SERVE\_FROM\_SUB\_PATH=true}).

% ============================================================
\section{Tests end-to-end et tests de charge}
% ============================================================

\subsection{Tests E2E --- Playwright}

Les tests end-to-end sont impl\'{e}ment\'{e}s avec \textbf{Playwright} au niveau
du Monorepo (\texttt{frontend/playwright.config.ts}). Trois profils de r\^{o}les
sont couverts~: administrateur, manager, technicien. Chaque profil poss\`{e}de
son propre fichier d'authentification sauvegard\'{e}
(\texttt{e2e/.auth/\{user,manager,technicien\}.json}), ce qui \'{e}vite une
reconnexion Keycloak \`{a} chaque sc\'{e}nario.

\subsubsection{Organisation des tests}

\begin{lstlisting}[language=bash, caption=Arborescence des tests E2E]
frontend/e2e/tests/
  auth.setup.ts              # Login admin       -> .auth/admin.json
  auth.setup.manager.ts      # Login manager     -> .auth/manager.json
  auth.setup.technicien.ts   # Login technicien  -> .auth/technicien.json
  dashboard.spec.ts          # Statistiques tableau de bord
  vehicles.spec.ts           # CRUD vehicules
  drivers.spec.ts            # CRUD conducteurs
  maintenance.spec.ts        # Suivi des interventions
  location.spec.ts           # Carte GPS temps reel
  roles.spec.ts              # Controle d'acces par role (RBAC)
\end{lstlisting}

\subsubsection{R\'{e}capitulatif des sc\'{e}narios}

\begin{table}[H]
\centering
\small
\renewcommand{\arraystretch}{1.45}
\begin{tabular}{
  >{\raggedright\arraybackslash}p{3cm}
  >{\raggedright\arraybackslash}p{8.5cm}
  >{\centering\arraybackslash}p{2cm}
}
\toprule
\rowcolor{secondary}
\textbf{Module} & \textbf{Sc\'{e}nario test\'{e}} & \textbf{R\'{e}sultat} \\
\midrule
Auth (admin)      & Connexion Keycloak + sauvegarde session        & \textcolor{success}{\textbf{Pass}} \\
\rowcolor{rowgray}
Auth (manager)    & Connexion avec r\^{o}le manager                & \textcolor{success}{\textbf{Pass}} \\
Auth (technicien) & Connexion avec r\^{o}le technicien             & \textcolor{success}{\textbf{Pass}} \\
\rowcolor{rowgray}
Dashboard         & Pr\'{e}sence des 4 cartes de statistiques      & \textcolor{success}{\textbf{Pass}} \\
Dashboard         & Valeurs charg\'{e}es sans erreur GraphQL        & \textcolor{success}{\textbf{Pass}} \\
\rowcolor{rowgray}
V\'{e}hicules     & Liste, cr\'{e}ation, suppression               & \textcolor{success}{\textbf{Pass}} \\
Conducteurs       & Liste et absence d'erreur GraphQL              & \textcolor{success}{\textbf{Pass}} \\
\rowcolor{rowgray}
Conducteurs       & Cr\'{e}ation (UUID v4, email unique)            & \textcolor{success}{\textbf{Pass}} \\
Conducteurs       & Modification (mise \`{a} jour t\'{e}l\'{e}phone) & \textcolor{success}{\textbf{Pass}} \\
\rowcolor{rowgray}
Conducteurs       & Changement de statut (\texttt{ON\_LEAVE})      & \textcolor{success}{\textbf{Pass}} \\
Conducteurs       & Suppression (soft delete v\'{e}rifi\'{e})       & \textcolor{success}{\textbf{Pass}} \\
\rowcolor{rowgray}
Maintenance       & Liste, planification, modification statut      & \textcolor{success}{\textbf{Pass}} \\
Localisation      & Chargement carte Leaflet                       & \textcolor{success}{\textbf{Pass}} \\
\rowcolor{rowgray}
RBAC              & Technicien redirig\'{e} depuis \texttt{/drivers}  & \textcolor{success}{\textbf{Pass}} \\
RBAC              & Manager redirig\'{e} depuis \texttt{/maintenance} & \textcolor{success}{\textbf{Pass}} \\
\bottomrule
\end{tabular}
\caption{R\'{e}capitulatif des 15 sc\'{e}narios end-to-end Playwright --- 0 \'{e}chec (Minikube)}
\end{table}

\begin{lstlisting}[language=bash, caption=Lancement des tests E2E]
./scripts/test-e2e.sh docker    # cible http://localhost:3005
./scripts/test-e2e.sh minikube  # cible http://flotte.local
\end{lstlisting}

\subsection{Tests de charge --- K6}

Les tests de charge sont impl\'{e}ment\'{e}s avec \textbf{K6} et couvrent les trois
protocoles utilis\'{e}s par la plateforme~: GraphQL (HTTP/JSON), REST et gRPC
(streaming).

\subsubsection{Profils de charge}

\begin{table}[H]
\centering
\renewcommand{\arraystretch}{1.45}
\begin{tabular}{
  >{\raggedright\arraybackslash}p{2.5cm}
  >{\centering\arraybackslash}p{2cm}
  >{\centering\arraybackslash}p{2.5cm}
  >{\raggedright\arraybackslash}p{5cm}
}
\toprule
\rowcolor{secondary}
\textbf{Profil} & \textbf{VUs max} & \textbf{Dur\'{e}e} & \textbf{Seuil d'acceptation} \\
\midrule
\textbf{Smoke}  & 1   & 1 min & p(99) $<$ 1\,500 ms \\
\rowcolor{rowgray}
\textbf{Load}   & 10  & 4 min & p(95) $<$ 500 ms \\
\textbf{Stress} & 100 & 4 min & taux d'erreur $<$ 1\,\% \\
\bottomrule
\end{tabular}
\caption{Profils de charge K6 d\'{e}finis dans \texttt{scenarios/options.js}}
\end{table}

Le profil \textit{load} suit un \textit{ramp-up} progressif~: 0 $\rightarrow$ 10 VUs
en 1 minute, maintien 2 minutes, descente en 1 minute, afin d'observer le
comportement de la plateforme sous charge r\'{e}aliste.

\subsubsection{Scripts de test}

\begin{lstlisting}[language=bash, caption=Arborescence des scripts K6]
load-tests/
  utils/config.js               # Auth Keycloak + helpers
  scenarios/options.js           # Profils smoke / load / stress
  targets/
    graphql-vehicles.js          # POST /graphql -> liste vehicules
    rest-alerts.js               # GET  /api/alerts -> liste alertes
    grpc-location.js             # gRPC WatchVehicle (server-streaming)
\end{lstlisting}

\begin{itemize}[leftmargin=*]
  \item \textbf{\texttt{graphql-vehicles.js}} : requ\^{e}te GraphQL
        \texttt{GetVehicles} avec jeton Bearer. V\'{e}rifie le statut HTTP 200,
        l'absence d'erreurs GraphQL et la pr\'{e}sence du tableau
        \texttt{vehicles}.
  \item \textbf{\texttt{rest-alerts.js}} : appel REST
        \texttt{GET /api/alerts} avec authentification. V\'{e}rifie le statut
        200 et que la r\'{e}ponse est un tableau.
  \item \textbf{\texttt{grpc-location.js}} : ouvre un flux
        \texttt{LocationService/WatchVehicle} en gRPC server-streaming,
        teste la connexion et le half-close propre.
\end{itemize}

\begin{lstlisting}[language=bash, caption=Lancement des tests K6]
./scripts/load-tests.sh docker    # smoke -> load -> stress (Docker Compose)
./scripts/load-tests.sh minikube  # smoke -> load -> stress (Minikube)
\end{lstlisting}

% ============================================================
\section{D\'{e}ploiement final avec Helm}
% ============================================================

\subsection{Architecture des charts}

L'ensemble du syst\`{e}me est d\'{e}ploy\'{e} via trois charts Helm ind\'{e}pendants,
permettant de mettre \`{a} jour chaque couche s\'{e}par\'{e}ment~:

\begin{lstlisting}[language=bash, caption=Structure des charts Helm]
infra/helm/
  fleet-infra/          # PostgreSQL (TimescaleDB), Kafka (KRaft), Redis
  fleet-observability/  # Loki, Grafana, Jaeger, Prometheus, OTel, Promtail
  fleet-app/            # 8 services applicatifs + frontend + ingress
\end{lstlisting}

\begin{itemize}[leftmargin=*]
  \item \textbf{\texttt{fleet-infra}} : infrastructure stateful. G\`{e}re les bases
        de donn\'{e}es PostgreSQL/TimescaleDB par service, Kafka en mode KRaft
        (sans ZooKeeper) et Redis pour le cache.
  \item \textbf{\texttt{fleet-observability}} : stack d'observabilit\'{e}. Toutes
        les ressources Kubernetes (Deployments, Services, ConfigMaps) sont
        g\'{e}n\'{e}r\'{e}es depuis les templates du chart.
  \item \textbf{\texttt{fleet-app}} : couche applicative. D\'{e}ploie les 6
        microservices m\'{e}tier, la Gateway GraphQL, le frontend React, Keycloak
        et les r\`{e}gles d'ingress.
\end{itemize}

\subsection{D\'{e}ploiement automatis\'{e} --- \texttt{scripts/kube.sh}}

Le script \texttt{kube.sh} orchestre l'int\'{e}gralit\'{e} du d\'{e}ploiement Minikube
en une seule commande~:

\begin{lstlisting}[language=bash, caption=\'{E}tapes du script \texttt{kube.sh}]
# 1. Demarrage Minikube + addon ingress
minikube start --cpus=4 --memory=8192
minikube addons enable ingress

# 2. Namespace et secrets
kubectl create namespace flotte-namespace
kubectl create secret generic db-secrets ...

# 3. Build des images dans le daemon Docker de Minikube
eval $(minikube docker-env)
docker build -t vehicle-service:latest ./services/vehicle-service
# ... (driver, maintenance, events, location, gateway, frontend)

# 4. Installation des charts dans l'ordre
helm install fleet-infra      ./infra/helm/fleet-infra
helm install fleet-observability ./infra/helm/fleet-observability

# 5. Creation des topics Kafka
kubectl exec -n flotte-namespace kafka-0 -- kafka-topics.sh \
  --create --topic flotte.location.events ...

# 6. Deploiement applicatif
helm install fleet-app ./infra/helm/fleet-app
\end{lstlisting}

\subsection{Points d'acc\`{e}s}

\begin{table}[H]
\centering
\renewcommand{\arraystretch}{1.4}
\begin{tabular}{
  >{\raggedright\arraybackslash}p{4cm}
  >{\ttfamily\raggedright\arraybackslash}p{7cm}
}
\toprule
\rowcolor{secondary}
\textbf{Service} & \textbf{URL} \\
\midrule
Frontend shell   & http://flotte.local \\
\rowcolor{rowgray}
GraphQL Gateway  & http://flotte.local/graphql \\
Keycloak         & http://flotte.local/auth \\
\rowcolor{rowgray}
Grafana          & http://flotte.local/grafana \\
Jaeger UI        & http://flotte.local/jaeger \\
\bottomrule
\end{tabular}
\caption{Points d'acc\`{e}s de la plateforme sous Minikube}
\end{table}

% ============================================================
\section{Page d'alertes dans le shell frontend}
% ============================================================

\subsection{Choix d'int\'{e}gration}

La page \texttt{/alerts} a \'{e}t\'{e} int\'{e}gr\'{e}e \textbf{directement dans l'application
shell} (\texttt{shell/src/alerts/AlertsPage.tsx}) plut\^{o}t que
comme micro-frontend d\'{e}di\'{e}. Ce choix se justifie par deux raisons~:

\begin{enumerate}[leftmargin=*]
  \item \textbf{Taille fonctionnelle} : la fonctionnalit\'{e} tient en un seul
        composant React~; cr\'{e}er un remote Module Federation pour cela aurait
        ajout\'{e} de la complexit\'{e} de build sans valeur ajout\'{e}e.
  \item \textbf{Transversalit\'{e}} : les alertes concernent tous les r\^{o}les et
        doivent \^{e}tre accessibles depuis n'importe quel contexte de navigation~;
        les placer dans le shell garantit leur disponibilit\'{e} permanente.
\end{enumerate}

\subsection{Fonctionnalit\'{e}s de \texttt{AlertsPage.tsx}}

\subsubsection{Temps r\'{e}el via subscription GraphQL}

Les nouvelles alertes apparaissent instantan\'{e}ment sans rechargement gr\^{a}ce
\`{a} une subscription WebSocket \texttt{ALERT\_CREATED\_SUB} via Apollo Client.

\begin{lstlisting}[language=Java, caption=Subscription temps r\'{e}el des alertes
--- \texttt{AlertsPage.tsx}]
const ALERT_CREATED_SUB = gql`
  subscription OnAlertCreated {
    alertCreated {
      id
      type
      severity
      status
      message
      vehicleId
      driverId
      createdAt
    }
  }
`;

useSubscription(ALERT_CREATED_SUB, {
  onData: ({ data }) => {
    const newAlert = data.data?.alertCreated;
    if (newAlert) {
      setAlerts(prev => [newAlert, ...prev]);
    }
  },
});
\end{lstlisting}

\subsubsection{Filtres et pagination}

\begin{itemize}[leftmargin=*]
  \item Filtre par \textbf{statut}~: ACTIVE, ACKNOWLEDGED, RESOLVED, EXPIRED
  \item Filtre par \textbf{s\'{e}v\'{e}rit\'{e}}~: CRITICAL, HIGH, WARNING, INFO
  \item Pagination~: 20 alertes par page
\end{itemize}

\subsubsection{Actions selon le r\^{o}le (RBAC)}

\begin{table}[H]
\centering
\renewcommand{\arraystretch}{1.4}
\begin{tabular}{
  >{\raggedright\arraybackslash}p{4.5cm}
  >{\centering\arraybackslash}p{2.5cm}
  >{\centering\arraybackslash}p{2.5cm}
  >{\centering\arraybackslash}p{2.5cm}
}
\toprule
\rowcolor{secondary}
\textbf{Action} & \textbf{admin} & \textbf{manager} & \textbf{technicien} \\
\midrule
Acquitter (\textit{acknowledge}) & \textcolor{success}{\checkmark} & \textcolor{success}{\checkmark} & \texttimes \\
\rowcolor{rowgray}
R\'{e}soudre (\textit{resolve})  & \textcolor{success}{\checkmark} & \textcolor{success}{\checkmark} & \textcolor{success}{\checkmark} \\
\bottomrule
\end{tabular}
\caption{Contr\^{o}le d'acc\`{e}s aux actions sur les alertes}
\end{table}

\subsubsection{Types d'alertes affich\'{e}s}

\begin{table}[H]
\centering
\small
\renewcommand{\arraystretch}{1.4}
\begin{tabular}{
  >{\ttfamily\raggedright\arraybackslash}p{5.8cm}
  >{\raggedright\arraybackslash}p{6.5cm}
}
\toprule
\rowcolor{secondary}
\normalfont\textbf{Code} & \normalfont\textbf{Label affich\'{e}} \\
\midrule
GEOFENCING\_BREACH      & Sortie de zone \\
\rowcolor{rowgray}
SPEED\_EXCEEDED         & Exc\`{e}s de vitesse \\
VEHICLE\_IMMOBILIZED    & Immobilisation v\'{e}hicule \\
\rowcolor{rowgray}
MAINTENANCE\_OVERDUE    & Maintenance en retard \\
LICENSE\_EXPIRING       & Permis expirant \\
\rowcolor{rowgray}
LICENSE\_EXPIRED        & Permis expir\'{e} \\
DELIVERY\_TOUR\_DELAYED & Tourn\'{e}e en retard \\
\bottomrule
\end{tabular}
\caption{Types d'alertes m\'{e}tier et leurs libell\'{e}s dans l'interface}
\end{table}

\subsection{Int\'{e}gration dans le Layout et la navigation}

La route \texttt{/alerts} est d\'{e}clar\'{e}e dans l'application Shell et ajout\'{e}e
au menu de navigation du \texttt{Layout.tsx}. Elle est visible par tous les
r\^{o}les (admin, manager, technicien).

\begin{lstlisting}[language=Java, caption=Route /alerts dans \texttt{App.tsx} (shell)]
import AlertsPage from './alerts/AlertsPage';

<Route path="alerts" element={<AlertsPage />} />
\end{lstlisting}

% ============================================================
\section{Pr\'{e}paration de la d\'{e}monstration}
% ============================================================

\subsection{Script de simulation d'alertes --- \texttt{simulate-alerts.sh}}

Le script \texttt{scripts/simulate-alerts.sh} g\'{e}n\`{e}re des sc\'{e}narios r\'{e}alistes
en s'appuyant sur les v\'{e}hicules et conducteurs r\'{e}els pr\'{e}sents en base (r\'{e}cup\'{e}r\'{e}s
via une requ\^{e}te GraphQL). Il publie ensuite les \'{e}v\'{e}nements directement
dans Kafka.

\begin{lstlisting}[language=bash, caption=Lancement de la simulation]
./scripts/simulate-alerts.sh minikube  # ou docker
\end{lstlisting}

\begin{table}[H]
\centering
\small
\renewcommand{\arraystretch}{1.4}
\begin{tabular}{
  >{\raggedright\arraybackslash}p{4.4cm}
  >{\ttfamily\raggedright\arraybackslash}p{5.5cm}
  >{\raggedright\arraybackslash}p{3.2cm}
}
\toprule
\rowcolor{secondary}
\normalfont\textbf{Type d'\'{e}v\'{e}nement} & \normalfont\textbf{Topic Kafka} & \normalfont\textbf{D\'{e}tail} \\
\midrule
SPEED\_EXCEEDED        & flotte.location.events    & 138, 108, 65 km/h \\
\rowcolor{rowgray}
GEOFENCING\_BREACH     & flotte.location.events    & Zone Nord, P\'{e}rim\`{e}tre Rouen \\
VEHICLE\_IMMOBILIZED   & flotte.location.events    & V\'{e}hicule arr\^{e}t\'{e} trop longtemps \\
\rowcolor{rowgray}
MAINTENANCE\_OVERDUE   & flotte.alertes.events     & S\'{e}v\'{e}rit\'{e} CRITICAL et HIGH \\
LICENSE\_EXPIRING      & flotte.conducteurs.events & 5 et 12 jours restants \\
\rowcolor{rowgray}
LICENSE\_EXPIRED       & flotte.conducteurs.events & Permis expir\'{e} \\
\bottomrule
\end{tabular}
\caption{\'{E}v\'{e}nements g\'{e}n\'{e}r\'{e}s par le script de simulation}
\end{table}

\subsection{Script de simulation GPS --- \texttt{simulate-GPS.sh}}

Le script \texttt{scripts/simulate-GPS.sh} simule un v\'{e}hicule en mouvement
en envoyant des coordonn\'{e}es GPS r\'{e}elles via gRPC directement vers le
\texttt{location-service}. Le trajet part de Rouen et progresse en temps
r\'{e}el~; la carte Leaflet du module \texttt{location} se met \`{a} jour
automatiquement via la subscription WebSocket.

\begin{lstlisting}[language=bash, caption=Lancement de la simulation GPS]
./scripts/simulate-GPS.sh minikube  # ouvre le port-forward gRPC si besoin
./scripts/simulate-GPS.sh docker
\end{lstlisting}

\subsection{Surveillance Kafka en temps r\'{e}el --- \texttt{watch-kafka.sh}}

Pendant la d\'{e}monstration, le script \texttt{scripts/watch-kafka.sh} permet
d'afficher en direct les \'{e}v\'{e}nements circulant sur tous les topics
\texttt{flotte.*}. Il supporte les deux environnements~:

\begin{lstlisting}[language=bash, caption=Surveillance Kafka en temps r\'{e}el]
./scripts/watch-kafka.sh docker    # via docker exec sur flotte-kafka
./scripts/watch-kafka.sh minikube  # via kubectl exec sur le pod Kafka
\end{lstlisting}

Cela permet de montrer le flux \'{e}v\'{e}nementiel complet~:
de la publication GPS jusqu'\`{a} l'apparition de l'alerte dans le navigateur.

\subsection{Collection Bruno --- Tests manuels API}

La collection Bruno (\texttt{bruno/}) permet de rejouer manuellement
n'importe quel appel API pendant la d\'{e}monstration~:

\begin{table}[H]
\centering
\renewcommand{\arraystretch}{1.4}
\begin{tabular}{
  >{\raggedright\arraybackslash}p{4cm}
  >{\centering\arraybackslash}p{2cm}
  >{\raggedright\arraybackslash}p{6cm}
}
\toprule
\rowcolor{secondary}
\textbf{Cat\'{e}gorie} & \textbf{Nb requ\^{e}tes} & \textbf{Services couverts} \\
\midrule
GraphQL   & 34 & V\'{e}hicules, conducteurs, maintenance, alertes, localisation \\
\rowcolor{rowgray}
REST      & 37 & CRUD complet de tous les microservices \\
gRPC      & 2  & WatchVehicle, StreamPositions \\
\bottomrule
\end{tabular}
\caption{Collection Bruno~: 73 requ\^{e}tes API pr\^{e}tes pour la d\'{e}monstration}
\end{table}

L'authentification Keycloak est g\'{e}r\'{e}e automatiquement par un
\textit{pre-request script} Bruno~: il \'{e}change les credentials contre un
JWT, d\'{e}tecte l'expiration (marge de 120 secondes) et injecte le token
dans l'en-t\^{e}te \texttt{Authorization} de chaque requ\^{e}te. Deux
environnements sont pr\'{e}-configur\'{e}s~: \texttt{docker} et \texttt{minikube}.

% ============================================================
\newpage
\section{Script d'installation --- \texttt{install.sh}}
% ============================================================

Un script d'installation unifi\'{e} \texttt{install.sh} a \'{e}t\'{e} r\'{e}alis\'{e} \`{a} la
racine du d\'{e}p\^{o}t afin de permettre \`{a} n'importe quel utilisateur de
d\'{e}marrer la plateforme compl\`{e}te en une seule commande, sans connaissance
pr\'{e}alable de l'architecture interne.

\begin{lstlisting}[language=bash, caption=Lancement de l'installation]
bash install.sh
\end{lstlisting}

\subsection{V\'{e}rification des pr\'{e}requis}

La premi\`{e}re \'{e}tape contr\^{o}le la pr\'{e}sence de tous les outils n\'{e}cessaires
avant toute action. En cas de manquant, le script liste les liens
d'installation officiels et s'arr\^{e}te imm\'{e}diatement.

\begin{table}[H]
\centering
\renewcommand{\arraystretch}{1.4}
\begin{tabular}{
  >{\ttfamily\raggedright\arraybackslash}p{4cm}
  >{\raggedright\arraybackslash}p{8.5cm}
}
\toprule
\rowcolor{secondary}
\normalfont\textbf{Outil} & \normalfont\textbf{Usage dans le projet} \\
\midrule
docker            & Build et ex\'{e}cution des conteneurs (mode Docker Compose) \\
\rowcolor{rowgray}
docker-compose    & Orchestration des services en d\'{e}veloppement \\
node / npm        & Installation des d\'{e}pendances workspaces et build frontend \\
\rowcolor{rowgray}
k6                & Ex\'{e}cution des tests de charge (smoke / load / stress) \\
minikube          & Cluster Kubernetes local (mode Minikube uniquement) \\
\rowcolor{rowgray}
helm              & D\'{e}ploiement des charts \texttt{fleet-*} \\
kubectl           & Gestion du cluster et cr\'{e}ation des topics Kafka \\
\bottomrule
\end{tabular}
\caption{Pr\'{e}requis v\'{e}rifi\'{e}s automatiquement par \texttt{install.sh}}
\end{table}

\subsection{Installation des d\'{e}pendances npm}

Apr\`{e}s la v\'{e}rification des pr\'{e}requis, le script ex\'{e}cute \texttt{npm install}
\`{a} la racine du monorepo. Gr\^{a}ce \`{a} la configuration npm workspaces, cette
commande installe en une seule passe les d\'{e}pendances des trois espaces de
travail (\texttt{gateway}, \texttt{services/location-service}, \texttt{frontend})
avec d\'{e}duplication automatique.

\subsection{S\'{e}lection interactive de l'environnement}

L'utilisateur choisit son environnement cible via un menu interactif~:

\begin{lstlisting}[language=bash, caption=Menu de s\'{e}lection de l'environnement]
  1) Docker Compose  (simple, recommande pour le developpement)
  2) Kubernetes / Minikube  (avance, production-like)
  3) Passer  (vous demarrerez manuellement)
\end{lstlisting}

\subsubsection{Mode Docker Compose}

Le script ex\'{e}cute \texttt{docker-compose up -d --build}, qui construit toutes
les images localement et d\'{e}marre l'ensemble des services. Une fois les
conteneurs actifs, les points d'acc\`{e}s sont affich\'{e}s~:

\begin{table}[H]
\centering
\renewcommand{\arraystretch}{1.4}
\begin{tabular}{
  >{\raggedright\arraybackslash}p{4.5cm}
  >{\ttfamily\raggedright\arraybackslash}p{6cm}
}
\toprule
\rowcolor{secondary}
\textbf{Service} & \textbf{URL} \\
\midrule
Frontend shell   & http://localhost:3005 \\
\rowcolor{rowgray}
GraphQL Gateway  & http://localhost:4000 \\
Keycloak         & http://localhost:8180 \\
\rowcolor{rowgray}
Grafana          & http://localhost:3001 \\
Jaeger UI        & http://localhost:16686 \\
\bottomrule
\end{tabular}
\caption{Points d'acc\`{e}s en mode Docker Compose}
\end{table}

\subsubsection{Mode Minikube}

En mode Minikube, le script d\'{e}l\`{e}gue enti\`{e}rement \`{a}
\texttt{scripts/kube.sh} (d\'{e}crit section~4.2). \`{A} l'issue du
d\'{e}ploiement, il affiche les instructions de configuration r\'{e}seau
sp\'{e}cifiques \`{a} la plateforme de l'utilisateur~:

\begin{itemize}[leftmargin=*]
  \item Ajout de \texttt{127.0.0.1 flotte.local} dans \texttt{/etc/hosts}.
  \item Sur WSL2 / Windows / macOS~: lancement de \texttt{minikube tunnel}
        dans un terminal s\'{e}par\'{e} pour exposer les services via l'ingress.
  \item Sur Linux natif~: aucune action suppl\'{e}mentaire n\'{e}cessaire.
\end{itemize}

\subsection{R\'{e}sum\'{e} final et prochaines \'{e}tapes}

Quel que soit l'environnement choisi, le script affiche un r\'{e}sum\'{e} des
prochaines \'{e}tapes~: URL de connexion, identifiants par d\'{e}faut
(\texttt{admin~/ admin}), et rappel des commandes utiles
(\texttt{docker-compose ps}, \texttt{docker-compose logs}, etc.).

\begin{infobox}
\textbf{Robustesse~:} le script utilise \texttt{set -e} pour s'arr\^{e}ter
imm\'{e}diatement en cas d'erreur, et d\'{e}tecte automatiquement la commande
disponible (\texttt{docker-compose} ou \texttt{docker compose}) pour assurer
la compatibilit\'{e} avec les versions r\'{e}centes de Docker.
\end{infobox}

% ============================================================
\newpage
\section{Refonte de l'architecture du d\'{e}pot}
% ============================================================

\subsection{Migration vers npm workspaces}

Un probl\`{e}me architectural majeur a \'{e}t\'{e} identifi\'{e} et corrig\'{e}~:
le d\'{e}pot contenait plusieurs \texttt{package.json} sans configuration de
monorepo, causant une duplication des d\'{e}pendances et une gestion des versions
incoh\'{e}rente.

\subsubsection{Structure pr\'{e}c\'{e}dente}

\begin{itemize}[leftmargin=*]
  \item \texttt{package.json} racine vide (seulement 150 octets, d\'{e}pendances
        gRPC incoh\'{e}rentes)
  \item \texttt{gateway/package.json} (Express + Apollo)
  \item \texttt{services/location-service/package.json} (NestJS + gRPC)
  \item \texttt{frontend/package.json} (React micro-frontend)
  \item 3 fichiers \texttt{package-lock.json} ind\'{e}pendants $\Rightarrow$ versions
        diff\'{e}rentes install\'{e}es en parall\`{e}le
\end{itemize}

\subsubsection{Nouvelle architecture (npm workspaces)}

Le \texttt{package.json} racine a \'{e}t\'{e} reconfigur\'{e} comme \textit{monorepo}
centralisant les d\'{e}pendances partag\'{e}es et la gestion des builds~:

\begin{lstlisting}[language=bash, caption=Configuration npm workspaces dans \texttt{package.json}]
{
  "name": "archi-dist-flotte-vehicules",
  "workspaces": [
    "services/location-service",
    "gateway",
    "frontend"
  ],
  "scripts": {
    "install": "npm install --workspaces",
    "build": "npm run build --workspaces",
    "dev": "npm run dev --workspaces",
    "test": "npm run test --workspaces",
    "start:gateway": "npm start --workspace=gateway",
    "start:location": "npm start --workspace=services/location-service"
  }
}
\end{lstlisting}

\subsubsection{B\'{e}n\'{e}fices imm\'{e}diats}

\begin{itemize}[leftmargin=*]
  \item \textbf{Deduplication automatique} : npm installe une seule copie de
        \texttt{@grpc/grpc-js} (v1.14.3) et la partage via symlinks.
  \item \textbf{Lock file centralis\'{e}} : un seul \texttt{package-lock.json} \`{a}
        la racine, suppression des fichiers en doublons (\texttt{gateway/} et
        \texttt{frontend/}).
  \item \textbf{Commandes unifi\'{e}es} : \texttt{npm run build --workspaces} compile
        tous les services d'un coup, simplifiant le CI/CD.
  \item \textbf{Versionning coh\'{e}rent} : tous les workspaces utilisent les m\^{e}mes
        versions de d\'{e}pendances, \'{e}vitant les incompatibilit\'{e}s subtiles.
\end{itemize}

\subsection{Nettoyage des r\'{e}sidus}

Les fichiers suivants ont \'{e}t\'{e} supprim\'{e}s pour \'{e}viter les conflits~:
\begin{itemize}[leftmargin=*]
  \item \texttt{gateway/package-lock.json}
  \item \texttt{frontend/package-lock.json}
\end{itemize}

Seul le \texttt{package-lock.json} racine demeure, maintenu par npm
automatiquement lors de chaque \texttt{npm install}.

\end{document}